\subsection{FIRST CONSEQUENCES}

C6. Let \(A, B, C\) be relations in \(\mathfrak{C}\) . If \(A \Longrightarrow B\) and \(B \Longrightarrow C\) are theorems in \(\mathfrak{C}\) , then \(A \Longrightarrow C\) is a theorem in \(\mathfrak{C}\) .

For \((\pmb{B} \Rightarrow \pmb{C}) \Rightarrow ((\pmb{A} \Rightarrow \pmb{B}) \Rightarrow (\pmb{A} \Rightarrow \pmb{C}))\) is an axiom of \(\mathfrak{C}\) , by replacing \(\pmb{A}\) by \(\pmb{B}\) , \(\pmb{B}\) by \(\pmb{C}\) , and \(\pmb{C}\) by ``not \(\pmb{A}\)'' in S4. By C1 (§2, no. 2), \((\pmb{A} \Rightarrow \pmb{B}) \Rightarrow (\pmb{A} \Rightarrow \pmb{C})\) is a theorem in \(\mathfrak{C}\) . A further application of C1 completes the proof.

C7. If A and B are relations in C, then \(\mathbf{B} \Longrightarrow (\mathbf{A} \text{ or } \mathbf{B})\) is a theorem in C.

For \(\boldsymbol{B}\Longrightarrow(\boldsymbol{B}\) or A) and \((\boldsymbol{B}\) or A) \(\Longrightarrow(\boldsymbol{A}\) or B) are axioms of C by virtue of S2 and S3. Now use C6.

C8. If \(A\) is a relation in \(\mathfrak{C}\) , \(A \Longrightarrow A\) is a theorem in \(\mathfrak{C}\) .

For \(A \Longrightarrow (A \text{ or } A)\) and \((A \text{ or } A) \Longrightarrow A\) are axioms, by S2 and S1. Now use C6.

C9. If A is a relation and B a theorem in T, then \(A \Longrightarrow B\) is a theorem in T.

For \(\boldsymbol{B} \Rightarrow ((\text{not } \boldsymbol{A}) \text{ or } \boldsymbol{B})\) is a theorem by C7, and therefore ``(not \(\boldsymbol{A}\) ) or \(\boldsymbol{B}\) '', that is to say \(\boldsymbol{A} \Rightarrow \boldsymbol{B}\) , is a theorem by C1.

C10. If A is a relation in T, then ``A or (not A)'' is a theorem in T.

For ``(not A) or A'' is a theorem by C8; now use S3 and C1.

C11. If A is a relation in C, ``A \(\Longrightarrow\) (not not A)'' is a theorem in C.

For this relation is ``(not A) or (not not A)'', and the result follows from C10.

C12. Let A and B be two relations in C. Then the relation

\[
(A \Rightarrow B) \Rightarrow ((\text { not } B) \Rightarrow (\text { not } A))
\]

is a theorem in \(\mathfrak{C}\) .

For

\[
((\text { not } A) \text { or } B) \Longrightarrow ((\text { not } A) \text { or } (\text { not   not } B))
\]

is a theorem, by C11, S4 and C1. On the other hand,

\[
((\text { not } A) \text { or } (\text { not   not } B)) \Longrightarrow ((\text { not   not } B) \text { or } (\text { not } A))
\]

is an axiom, by S3. Therefore

\[
((\text { not } A) \text { or } B) \Longrightarrow ((\text { not   not } B) \text { or } (\text { not } A))
\]

is a theorem by C6. Hence the result.

C13. Let \(A, B, C\) be relations in \(\mathfrak{C}\) . If \(A \Rightarrow B\) is a theorem in \(\mathfrak{C}\) , then \((B \Rightarrow C) \Rightarrow (A \Rightarrow C)\) is a theorem in \(\mathfrak{C}\) .

For (not B) \(\Longrightarrow\) (not A) is a theorem, by C12 and C1. Therefore (C or (not B)) \(\Longrightarrow\) (C or (not A)) is a theorem, by S4 and C1. By a double application of S3 and C6 we infer that

\[
((\text { not } \boldsymbol {B}) \text { or } \boldsymbol {C}) \Longrightarrow ((\text { not } \boldsymbol {A}) \text { or } \boldsymbol {C})
\]

is a theorem; but this is the given relation.

¶ From now on, we shall generally use C1 and C6 without quoting them explicitly.

